\lecture{Lecture 13: tbd}

\qa{How does PCA-SIFT alter the standard SIFT descriptor calculation, and why does the patch size strictly shrink from \texorpdfstring{$41 \times 41$}{41x41} to \texorpdfstring{$39 \times 39$}{39x39}?}
{
\begin{itemize}
    \item PCA-SIFT retains the first three steps of SIFT (detection, scale, orientation) but fundamentally changes the descriptor extraction. It extracts a $41 \times 41$ pixel patch centered on the keypoint and rotated to the canonical orientation.
    \item Instead of building 8-bin orientation histograms, it computes the raw horizontal and vertical gradients for the patch. 
    \item The bounding pixels are strictly discarded (reducing the patch to $39 \times 39$) because calculating a discrete spatial derivative requires comparing a pixel to its neighbors, which is mathematically impossible for pixels on the outermost edge.
    \item The horizontal and vertical gradients are concatenated into a massive vector ($2 \times 39 \times 39 = 3042$ dimensions). This vector is normalized and then projected onto a lower-dimensional linear space using \textbf{Principal Component Analysis (PCA)} to maximize variance while drastically reducing the descriptor's footprint.
\end{itemize}}

\qa{How does SURF utilize integral images and box filters to achieve its significant speed advantage over SIFT's Hessian matrix computation?}
{
\begin{itemize}
    \item \textbf{Integral Image:} An intermediate representation $I_\Sigma(x,y)$ where each pixel value is the sum of all pixels in the rectangle defined by the origin and $(x,y)$. This allows the sum of intensities inside \textit{any} arbitrary rectangular image area to be computed in $O(1)$ time using just four array references (e.g., $A + D - B - C$).
    \item \textbf{Box Filters:} SURF approximates the computationally heavy second-order Gaussian derivatives (for the Hessian matrix) using simple rectangular "box filters" composed of discrete black (negative weight) and white (positive weight) rectangles.
    \item Because these box filters are purely rectangular sums, they are computed instantly using the integral image, completely decoupling the processing time from the actual size of the filter.
\end{itemize}}

\qa{What is the structural difference in how the scale space pyramid is constructed in SURF compared to SIFT?}
{
\begin{itemize}
    \item \textbf{SIFT:} Iteratively smooths the image with a Gaussian filter and then physically downsamples (reduces the resolution of) the image to build the next octave.
    \item \textbf{SURF:} Never subsamples the original image. Instead of shrinking the image, SURF \textbf{up-scales the filter size}. It applies progressively larger box filters directly to the original-resolution integral image. Because integral image lookups are $O(1)$, applying a massive box filter costs the exact same computational time as a small one, eliminating the need for iterative image resizing.
\end{itemize}}

\qa{How does SURF determine the dominant keypoint orientation without computing exact, continuous image gradients?}
{
\begin{itemize}
    \item SURF evaluates a circular neighborhood of radius $6s$ ($s$ = scale). 
    \item It approximates the $x$ and $y$ gradients extremely fast using \textbf{Haar wavelets} (simple binary black/white filters).
    \item The responses are plotted on a 2D $(x,y)$ plane. A sliding wedge of $60^\circ$ sweeps around the origin, summing the wavelet responses within it. The angle of the wedge that yields the highest total sum is assigned as the dominant keypoint orientation.
\end{itemize}}

\qa{What is the mathematical composition of the final SURF descriptor, and what is its primary engineering trade-off compared to SIFT?}
{
\begin{itemize}
    \item \textbf{Composition:} SURF extracts a square region around the keypoint (sized $20s \times 20s$, where $s$ is the scale) aligned with the dominant orientation, and divides it into a $4 \times 4$ spatial grid of sub-regions (16 total). For each sub-region, instead of 8-bin histograms, SURF computes a simple 4-dimensional vector using Haar wavelet responses: $v = (\sum d_x, \sum d_y, \sum |d_x|, \sum |d_y|)$. 
    \item Concatenating these 16 sub-regions yields a $16 \times 4 = \textbf{64-dimensional descriptor}$ (exactly half the size of standard SIFT).
    \item \textbf{Trade-off:} SURF achieves drastically faster computation and matching times due to integral images and a smaller descriptor vector. However, it is structurally \textbf{less robust} to extreme illumination variations and viewpoint (perspective) changes compared to SIFT.
\end{itemize}}

\qa{How does the Gradient Location-Orientation Histogram (GLOH) alter SIFT's spatial grid, and how does it manage the resulting dimensionality explosion?}
{
\begin{itemize}
    \item GLOH modifies the descriptor calculation step by replacing SIFT's Cartesian $4 \times 4$ grid with a \textbf{log-polar location grid}. This grid consists of 3 radial distance bins and 8 angular bins (leaving the central base bin un-subdivided, resulting in $1 + 8 + 8 = 17$ spatial bins).
    \item Like SIFT, it computes 16 gradient orientations per spatial sub-region. This geometry creates $17 \times 16 = 272$ total orientation bins (a massive vector).
    \item To remain computationally competitive and remove redundant correlations, GLOH applies \textbf{Principal Component Analysis (PCA)} to strictly project the 272-dimensional vector down to a 128-dimensional space, matching SIFT's descriptor size but capturing radial spatial layouts more effectively.
\end{itemize}}

\qa{What is the mathematical basis of the Shape Context descriptor, and what critical limitation led to its obsolescence?}
{
\begin{itemize}
    \item It quantizes relative point locations into a \textbf{2D} log-polar coordinate system (5 radial bins, 12 angular bins). \textbf{For clarity}: while the spatial layout is 2D, the descriptor builds a \textbf{3D} histogram by additionally accumulating the \textbf{edge orientations} (e.g., extracted via Canny) as the third dimension, counting the discrete number of edge points falling into each bin.
    \item Scale invariance is achieved by normalizing all radial distances by the mean distance between points.
    \item \textbf{Fatal Flaw:} While robust to noise and minor deformations, it is strictly \textbf{dependent on rotations}. It lacks native rotation invariance, making it highly brittle to viewpoint changes compared to modern descriptors.
\end{itemize}}

\qa{How does LBP encode local image patches, and what is its primary architectural advantage?}
{
\begin{itemize}
    \item LBP compares a central pixel's intensity to its 8 immediate neighbors. Using a thresholding function: if the neighbor is strictly greater than the center, it outputs 1; otherwise, 0.
    \item This generates an 8-bit binary sequence (converted to a simple integer). Histograms of these values are concatenated across sub-regions.
    \item \textbf{Advantage:} It is incredibly \textbf{compact} and does not strictly require a standalone keypoint detector. It can be applied densely across an entire image or bounding box (originally designed for robust text and texture recognition).
\end{itemize}}

\qa{How does BRIEF structurally differ from gradient-histogram descriptors like SIFT, and why is it vastly faster to match?}
{
\begin{itemize}
    \item Instead of accumulating continuous gradient vectors, BRIEF (Binary Robust Independent Elementary Features) applies Gaussian smoothing and directly compares the raw intensities of randomly sampled pixel pairs $(x,y)$ inside a window.
    \item It utilizes a binary test: $\tau(p; x,y) = 1$ if $p(x) < p(y)$, otherwise $0$. These binary results are concatenated into a single bitstring.
    \item \textbf{Speed Trade-off:} Because the descriptor is purely binary, feature matching relies on the \textbf{Hamming distance} (counting differing bits via a hardware-level XOR operation). This is computationally exponentially faster than calculating the Euclidean distance of floating-point vectors required by SIFT/SURF.
\end{itemize}}

\qa{What components make up the ORB algorithm, and how does it geometrically fix the main limitation of the BRIEF descriptor?}
{
\begin{itemize}
    \item ORB (Oriented FAST and Rotated BRIEF) fuses the highly efficient \textbf{FAST corner detector} with the \textbf{BRIEF descriptor}.
    \item Standard BRIEF is entirely rotation-dependent; if the image rotates, the fixed pixel-pair sampling fails.
    \item ORB fixes this by computing the \textbf{intensity centroid} of the keypoint neighborhood. The vector from the geometric center to the intensity centroid defines a dominant orientation. ORB then explicitly rotates the BRIEF pixel-pair sampling coordinates to align with this axis, creating a rotation-invariant binary descriptor.
\end{itemize}}

\qa{In feature space matching strategies, why is Nearest Neighbor Distance Ratio (NNDR) generally superior to simple Nearest Neighbor (NN)?}
{
\begin{itemize}
    \item \textbf{Simple NN} finds the mathematically closest descriptor in feature space but requires a hard, arbitrary global distance threshold to reject bad matches.
    \item \textbf{NNDR} evaluates the ratio between the distance to the absolute best match ($d_1$) and the second-best match ($d_2$). 
    \item If $d_1 \approx d_2$, it indicates an ambiguous feature (e.g., repeating patterns like windows on a building). NNDR dynamically filters out these structurally ambiguous matches without relying on a rigid distance threshold, keeping only highly distinct matches where $d_1 \ll d_2$.
\end{itemize}}

\qa{How do Precision and Recall mathematically and conceptually evaluate the quality of a feature matching pipeline?}
{
\begin{itemize}
    \item \textbf{Precision} ($TP / (TP + FP)$) measures \textit{exactness}: out of all the matches the algorithm claimed were correct, what fraction were genuinely correct? It heavily penalizes False Positives (wrongly matched features).
    \item \textbf{Recall} ($TP / (TP + FN)$) measures \textit{completeness}: out of all the true physical correspondences that actually exist in the scene, what fraction did the algorithm successfully detect? It penalizes False Negatives (missed true matches).
\end{itemize}}